% Generated draft: requires mathematical and source review.
\documentclass[11pt]{article}
\usepackage[margin=25mm]{geometry}
\usepackage{fontspec}
\setmainfont{Latin Modern Roman}
\usepackage{amsmath,amssymb,mathtools,unicode-math}
\setmathfont{Latin Modern Math}
\usepackage{enumitem,needspace,xurl,hyperref}
\setlength{\parindent}{0pt}
\setlength{\parskip}{5pt}
\setlength{\emergencystretch}{3em}
\newcommand{\N}{\mathbb{N}}
\newcommand{\Z}{\mathbb{Z}}
\newcommand{\Q}{\mathbb{Q}}
\newcommand{\R}{\mathbb{R}}
\newcommand{\C}{\mathbb{C}}
\newcommand{\F}{\mathbb{F}}
\newcommand{\A}{\mathbb{A}}
\newcommand{\PP}{\mathbb{P}}
\newcommand{\E}{\mathbb{E}}
\newcommand{\eps}{\varepsilon}
\newcommand{\dd}{\,\mathrm{d}}
\begin{document}
\section*{UnsolvedMath ID: 801051}
\textbf{Source problem code:} catalogue entry
\subsection{Definitions and mathematical statement}
Let \(\mathcal H\) denote the intended topological parameter space of quadratic H\'enon maps, let \(X\) denote their common Euclidean phase space, and let \(\lambda_X\) denote Lebesgue measure on that space. The supplied record does not specify the real or complex setting, the parameter domain, or its topology; these choices remain unresolved.

For a map \(H\colon X\to X\), a point \(x\in X\), and an integer \(n\geq 1\), define the empirical probability measure
\[
\mu_{n,H,x}=\frac1n\sum_{j=0}^{n-1}\delta_{H^j(x)},
\]
where \(H^0\) is the identity and \(\delta_y\) is the unit point mass at \(y\). Under the convention that non-statistical behavior means failure of these measures to converge weakly to a Borel probability measure on \(X\), put
\[
\mathcal N(H)=\left\{x\in X:\text{there is no }\mu\in\mathcal P(X)\text{ such that }\mu_{n,H,x}\Rightarrow\mu\text{ as }n\to\infty\right\}.
\]
Here \(\mathcal P(X)\) is the set of Borel probability measures on \(X\), and \(\mu_n\Rightarrow\mu\) means that \(\int_X\varphi\,\mathrm d\mu_n\to\int_X\varphi\,\mathrm d\mu\) for every bounded continuous function \(\varphi\colon X\to\R\). This convention is an explicit provisional interpretation: the supplied record does not define non-statistical behavior or specify how escaping orbits are treated.

\textbf{Question 3 (Berger and Bedford, 2023).} With the intended family, topology, phase space, and interpretation of non-statistical behavior fixed, do there exist \(\varepsilon>0\) and \(E\subseteq\mathcal H\) such that
\[
\operatorname{int}_{\mathcal H}\!\left(\overline E^{\,\mathcal H}\right)\neq\varnothing
\]
and
\[
\forall H\in E\quad\exists A_H\subseteq X\text{ Lebesgue measurable}:\qquad
\lambda_X(A_H)\geq\varepsilon
\quad\text{and}\quad
A_H\subseteq\mathcal N(H)?
\]
Thus the positive lower bound \(\varepsilon\) is common to all maps in \(E\), while the set \(A_H\) may depend on \(H\). The displayed closure condition is precisely the supplied meaning of locally dense. A definitive formulation requires resolving the choices identified above.
\subsection{Short English statement}
The question asks whether a locally dense collection of quadratic H\'enon maps can share a positive lower bound on the Lebesgue measure of a set consisting entirely of points with non-statistical behavior. The phase space, parameter topology, and meaning of non-statistical behavior require confirmation from the source; the empirical-measure formulation above is conditional on the stated convention.
\subsection{Sources}
\begin{itemize}
\item Local UnsolvedMath registry, permanent ID \(801051\), exact source problem code \texttt{catalogue entry}, title \emph{Henon Maps: Open Problems --- Question 3}. The record attributes the question to Berger and Bedford, with proposed year 2023, and supplies the status \texttt{open}; that status is not independently verified here.
\item Berger and Bedford, \emph{H\'enon Maps: A List of Open Problems} (2023), Question 3. Supplied source URL: \url{https://arxiv.org/abs/2312.03907}. The supplied provenance records access on 2026-10-16 and extraction by \texttt{arxiv\_latex}; these are supplied provenance statements, not a claim of access or verification in this draft.
\end{itemize}
\end{document}
